\NativeSlot{54}{36}{%
% p014_s01

\ExamQuestion{7（续）}{5}%
是\mbox{（\hspace{1.8em}）}。\par
}
\NativeSlot{72}{43}{%
\ExamGroup{二、选择（将正确答案的序号填在括号里）。{\normalfont\kaishu\fontsize{10.8bp}{14bp}\selectfont （每空2分，共16分）}}

% p014_s02

\ExamQuestion{1.}{5}%
在下面的字母中，线段的位置关系既有平行又有垂直的是\mbox{（\hspace{1.8em}）}。\par
\BookFigure{figures/p014_s02-01.tex}
}
\NativeSlot{147}{44}{%
% p014_s03

\ExamQuestion{2.}{5}%
我国战国时期《墨经》中记载“平，同高也”，它的意思是两条平行线间的距离\mbox{（\hspace{1.8em}）}。\par
\ExamOptions{A. 不相等}{B. 处处相等}{C. 有时相等}{D. 无法确定}
}
\NativeSlot{222}{44}{%
% p014_s04

\ExamQuestion{3.}{5}%
在同一平面内，三条直线两两相交，最多有\mbox{（\hspace{1.8em}）}个交点，最少有\mbox{（\hspace{1.8em}）}个交点。\par
\ExamOptions{A. 0}{B. 1}{C. 2}{D. 3}
}
\NativeSlot{297}{45}{%
% p014_s05

\ExamQuestion{4.}{5}%
若\(a\parallel b\)，\allowbreak{}\(b\parallel c\)，则\(a\)与\(c\)的位置关系是\mbox{（\hspace{6em}）}。\par
\ExamOptions{A. 互相平行}{B. 相交}{C. 互相垂直}{D. 可能互相平行，也可能互相垂直}
}
\NativeSlot{373}{59}{%
% p014_s06

\ExamQuestion{5.}{5}%
美术手工课上，同学们将一张正方形纸片对折2次，展开后得到的两条折痕\mbox{（\hspace{1.8em}）}。\par
\ExamOptions{A. 互相垂直}{B. 可能互相垂直，也可能互相平行}{C. 互相平行}{D. 既不互相垂直，也不互相平行}
}
\NativeSlot{472}{90}{%
% p014_s07

\ExamQuestion{6.}{5}%
某景区救援船停在人工湖的\(B\)点，发现\(A\)点有游客落水，救援船从\(B\)点赶往\(A\)点的最短行驶路线，对应的数学原理是\mbox{（\hspace{1.8em}）}。\par
\ExamOptions{A. 垂直线段最短}{B. 两点之间线段最短}{C. 平行线间距离相等}{D. 过一点有且只有一条直线与已知直线平行}
\BookFigure{figures/p014_s07-01.tex}
}
\NativeSlot{623}{80}{%
% p014_s08

\ExamQuestion{7.}{5}%
国家推进城乡供水一体化改造，在农村饮水安全工程中，从村庄垂直向河边拉输水管道，水管最短，节约管材、降低成本，其中蕴含的数学道理是\mbox{（\hspace{1.8em}）}。\par
\ExamOptions{A. 两点之间线段最短}{B. 垂直线段最短}{C. 两点确定一条直线}{D. 平行线间距离相等}
}
